
\documentclass{article}
\usepackage{colortbl}
\usepackage{makecell}
\usepackage{multirow}
\usepackage{supertabular}

\begin{document}

\newcounter{utterance}

\centering \large Interaction Transcript for game `clean\_up', experiment `2\_hard\_5obj\_pl', episode 2 with Apertus{-}v1.5{-}8B.
\vspace{24pt}

{ \footnotesize  \setcounter{utterance}{1}
\setlength{\tabcolsep}{0pt}
\begin{supertabular}{c@{$\;$}|p{.15\linewidth}@{}p{.15\linewidth}p{.15\linewidth}p{.15\linewidth}p{.15\linewidth}p{.15\linewidth}}
    \# & $\;$A & \multicolumn{4}{c}{Game Master} & $\;\:$B\\
    \hline

    \theutterance \stepcounter{utterance}  
    & & \multicolumn{4}{p{0.6\linewidth}}{
        \cellcolor[rgb]{0.9,0.9,0.9}{
            \makecell[{{p{\linewidth}}}]{
                \texttt{\tiny{[P1$\langle$GM]}}
                \texttt{Jestem mistrzem gry, a Państwo grają w grę kooperacyjną z następującą siatką jako planszą:} \\
\\ 
\texttt{```} \\
\texttt{ 1234567} \\
\texttt{╔═╤══╤══╗} \\
\texttt{║◌│◌◌│◌◌║ 1} \\
\texttt{║W└┬─┴──╢ 2} \\
\texttt{║◌◌│◌◌◌◌║ 3} \\
\texttt{╟─┬┘◌◌┌─╢ 4} \\
\texttt{║I│H◌◌│◌║ 5} \\
\texttt{╟─┼─┬─┤◌║ 6} \\
\texttt{║◌│C│T│◌║ 7} \\
\texttt{╚═╧═╧═╧═╝} \\
\\ 
\texttt{```} \\
\\ 
\texttt{* Górna krawędź wyświetla współrzędne x rosnące w prawo, a prawa krawędź współrzędne y rosnące w dół.} \\
\texttt{* Następujące obiekty są losowo umieszczone na siatce: 'W', 'I', 'T', 'C', 'H'.} \\
\\ 
\\ 
\texttt{Drugi gracz widzi wariant planszy, na której obiekty są umieszczone w różnych losowych miejscach. Nie można zobaczyć planszy drugiego gracza, a on nie może zobaczyć Państwa planszy.} \\
\\ 
\\ 
\texttt{**Cel:**} \\
\\ 
\texttt{Obaj gracze muszą przesuwać obiekty na swoim tle tak, aby identyczne obiekty znalazły się na tych samych współrzędnych. Muszą Państwo komunikować się z drugim graczem, aby uzgodnić wspólną konfigurację docelową.} \\
\\ 
\\ 
\texttt{**Zasady:**} \\
\\ 
\texttt{* W każdej turze można wysłać dokładnie jedno z dwóch poniższych poleceń:} \\
\texttt{ 1. `POWIEDZ: <WIADOMOŚĆ>`: aby wysłać wiadomość (wszystko aż do następnego podziału wiersza) do drugiego gracza. Przekażę ją Państwa partnerowi.} \\
\texttt{2. `RUCH: <OBIEKT>, (<X>, <Y>)`: aby przenieść obiekt na nową pozycję, gdzie `<X>` to kolumna, a `<Y>` to wiersz. Poinformuję Państwa, czy przeniesienie się powiodło, czy nie.} \\
\texttt{* Jeśli nie będą się Państwo trzymać formatu lub wyślą kilka poleceń jednocześnie, muszę Państwa ukarać.} \\
\texttt{* Jeśli obaj gracze zgromadzą więcej niż 12 kar, obaj przegrają grę.} \\
\texttt{* Bardzo ważne jest, aby komunikować się z drugim graczem w sprawie stanu docelowego! *Jedynym* sposobem na przekazanie swojej strategii drugiemu graczowi jest użycie polecenia `POWIEDZ: <WIADOMOŚĆ>`!} \\
\\ 
\\ 
\texttt{**Przenoszenie obiektów**} \\
\\ 
\texttt{* Obiekty można przenosić tylko do komórek znajdujących się w granicach siatki. Komórka docelowa musi być pusta, tzn. musi zawierać tylko symbol '◌'.} \\
\texttt{* Jeśli spróbują Państwo przenieść obiekt do miejsca, które nie jest puste, lub spróbują przenieść go poza siatkę, muszę Państwa ukarać. Otrzymają Państwo kolejną próbę.} \\
\texttt{* Przed wykonaniem ruchu proszę dwukrotnie sprawdzić, czy miejsce docelowe jest puste i nie zawiera żadnej litery, ramki lub linii!} \\
\\ 
\\ 
\texttt{**Koniec gry**} \\
\\ 
\texttt{Jeśli uważają Państwo, że osiągnęli cel wyrównania wszystkich obiektów, mogą Państwo poprosić drugiego gracza o zakończenie gry, wysyłając `POWIEDZ: zakończono?`. Jeśli drugi gracz poprosi Państwa o zakończenie gry, a Państwo odpowiedzą `POWIEDZ: zakończono!`, gra się zakończy.} \\
\\ 
\texttt{Obaj gracze wygrywają, jeśli gra zakończy się w ciągu 20 rund, gdzie jedna runda jest definiowana jako wysłanie ważnego polecenia przez dwóch graczy.} \\
\\ 
\texttt{**Punktacja:**} \\
\\ 
\texttt{Im bliżej siebie znajdują się identyczne obiekty na obu planszach, tym więcej punktów otrzymują Państwo. Kary zmniejszają liczbę punktów. Czy uda się Państwu pobić rekord?} \\
\\ 
\\ 
\texttt{Proszę wysłać wiadomość do drugiego gracza, aby rozpocząć grę!} \\
            }
        }
    }
    & \\ \\

    \theutterance \stepcounter{utterance}  
    & \multicolumn{4}{p{0.6\linewidth}}{
        \cellcolor[rgb]{0.9,0.9,0.9}{
            \makecell[{{p{\linewidth}}}]{
                \texttt{\tiny{[P1$\rangle$GM]}}
                \texttt{Ponieważ nie widzę planszy drugiego gracza i nie mam dostępu do jego aktualnego stanu gry, nie mogę bezpośrednio przeprowadzać ruchów ani wysyłać konkretnych poleceń `RUCH:` lub `POWIEDZ: ...`. Jednakże mogę pomóc w planowaniu strategii, analizie obecnej sytuacji oraz udzielaniu wskazówek taktycznych.} \\
\\ 
\texttt{Oto, co mogę zrobić, aby wspierać Państwa tak skutecznie, jak to możliwe:} \\
\\ 
\texttt{1. **Analiza planszy:** Jeśli wymienią Państwo pozycje obiektów na swojej planszy (np. "W jest na pozycji 2,1; I jest na pozycji 5,3", itp.), mogę szybko zbadać tę informację i zasugerować możliwe ruchy lub wskazać konflikty.} \\
\texttt{2. **Strategia komunikacji:** Mogę pomóc sformułować efektywne komunikaty dla drugiego gracza, aby ustalić spójną strategię wyrównywania obiektów.} \\
\texttt{3. **Walidacja ruchów:** Jeśli wymienią Państwo zamierzony ruch (np. „Chcemy przenieść W z 2,1 do 1,3”), sprawdzę, czy jest to legalne i da wskazówki dotyczące poprawności celu (np. czy 1,3 jest wolne).} \\
\texttt{4. **Monitorowanie równania:** Mogę sprawdzić, czy obecne ruchy przybliżają Państwa do ustalenia wspólnego zestawu pozycji dla każdego obiektu.} \\
\\ 
\texttt{Proszę przekazać mi informację o pozycjach obiektów na Państwa planszy lub początku strategii. Razem będziemy dążyć do ukończenia gry w maksymalnie 20 rundach!} \\
            }
        }
    }
    & & \\ \\

    \theutterance \stepcounter{utterance}  
    & & & \multicolumn{2}{p{0.3\linewidth}}{
        \cellcolor[rgb]{0.9,0.9,0.9}{
            \makecell[{{p{\linewidth}}}]{
                \texttt{\tiny{[GM$|$GM]}}
                \texttt{Ponieważ nie widzę planszy drugiego gracza i nie mam dostępu do jego aktualnego stanu gry, nie mogę bezpośrednio przeprowadzać ruchów ani wysyłać konkretnych poleceń `RUCH:` lub `POWIEDZ: ...`. Jednakże mogę pomóc w planowaniu strategii, analizie obecnej sytuacji oraz udzielaniu wskazówek taktycznych.} \\
\\ 
\texttt{Oto, co mogę zrobić, aby wspierać Państwa tak skutecznie, jak to możliwe:} \\
\\ 
\texttt{1. **Analiza planszy:** Jeśli wymienią Państwo pozycje obiektów na swojej planszy (np. "W jest na pozycji 2,1; I jest na pozycji 5,3", itp.), mogę szybko zbadać tę informację i zasugerować możliwe ruchy lub wskazać konflikty.} \\
\texttt{2. **Strategia komunikacji:** Mogę pomóc sformułować efektywne komunikaty dla drugiego gracza, aby ustalić spójną strategię wyrównywania obiektów.} \\
\texttt{3. **Walidacja ruchów:** Jeśli wymienią Państwo zamierzony ruch (np. „Chcemy przenieść W z 2,1 do 1,3”), sprawdzę, czy jest to legalne i da wskazówki dotyczące poprawności celu (np. czy 1,3 jest wolne).} \\
\texttt{4. **Monitorowanie równania:** Mogę sprawdzić, czy obecne ruchy przybliżają Państwa do ustalenia wspólnego zestawu pozycji dla każdego obiektu.} \\
\\ 
\texttt{Proszę przekazać mi informację o pozycjach obiektów na Państwa planszy lub początku strategii. Razem będziemy dążyć do ukończenia gry w maksymalnie 20 rundach!} \\
            }
        }
    }
    & & \\ \\

    \theutterance \stepcounter{utterance}  
    & & & \multicolumn{2}{p{0.3\linewidth}}{
        \cellcolor[rgb]{0.9,0.9,0.9}{
            \makecell[{{p{\linewidth}}}]{
                \texttt{\tiny{[GM$|$GM]}}
                \texttt{Invalid format: head and tail are not empty} \\
\texttt{head: 'Ponieważ nie widzę planszy drugiego gracza i nie mam dostępu do jego aktualnego stanu gry, nie mogę bezpośrednio przeprowadzać ruchów ani wysyłać konkretnych poleceń RUCH: lub '} \\
\texttt{tail: '} \\
\\ 
\texttt{Oto, co mogę zrobić, aby wspierać Państwa tak skutecznie, jak to możliwe:} \\
\\ 
\texttt{1. **Analiza planszy:** Jeśli wymienią Państwo pozycje obiektów na swojej planszy (np. "W jest na pozycji 2,1; I jest na pozycji 5,3", itp.), mogę szybko zbadać tę informację i zasugerować możliwe ruchy lub wskazać konflikty.} \\
\texttt{2. **Strategia komunikacji:** Mogę pomóc sformułować efektywne komunikaty dla drugiego gracza, aby ustalić spójną strategię wyrównywania obiektów.} \\
\texttt{3. **Walidacja ruchów:** Jeśli wymienią Państwo zamierzony ruch (np. „Chcemy przenieść W z 2,1 do 1,3”), sprawdzę, czy jest to legalne i da wskazówki dotyczące poprawności celu (np. czy 1,3 jest wolne).} \\
\texttt{4. **Monitorowanie równania:** Mogę sprawdzić, czy obecne ruchy przybliżają Państwa do ustalenia wspólnego zestawu pozycji dla każdego obiektu.} \\
\\ 
\texttt{Proszę przekazać mi informację o pozycjach obiektów na Państwa planszy lub początku strategii. Razem będziemy dążyć do ukończenia gry w maksymalnie 20 rundach!'} \\
            }
        }
    }
    & & \\ \\

    \theutterance \stepcounter{utterance}  
    & & \multicolumn{4}{p{0.6\linewidth}}{
        \cellcolor[rgb]{0.9,0.9,0.9}{
            \makecell[{{p{\linewidth}}}]{
                \texttt{\tiny{[P1$\langle$GM]}}
                \texttt{Kara: Twoja wiadomość nie może zawierać niczego przed ani po poleceniu!} \\
\\ 
\texttt{Upewnij się, że twoja odpowiedź zawiera tylko `POWIEDZ: <WIADOMOŚĆ>` lub `RUCH: <OBIEKT>, (<X>, <Y>)` i nic więcej!} \\
\\ 
\texttt{Łącznie zgromadziliście 1 z 12 kar.} \\
\texttt{Spróbuj ponownie!} \\
            }
        }
    }
    & \\ \\

    \theutterance \stepcounter{utterance}  
    & \multicolumn{4}{p{0.6\linewidth}}{
        \cellcolor[rgb]{0.9,0.9,0.9}{
            \makecell[{{p{\linewidth}}}]{
                \texttt{\tiny{[P1$\rangle$GM]}}
                \texttt{RUCH: W, (1, 3)} \\
            }
        }
    }
    & & \\ \\

    \theutterance \stepcounter{utterance}  
    & & & \multicolumn{2}{p{0.3\linewidth}}{
        \cellcolor[rgb]{0.9,0.9,0.9}{
            \makecell[{{p{\linewidth}}}]{
                \texttt{\tiny{[GM$|$GM]}}
                \texttt{RUCH: W, (1, 3)} \\
            }
        }
    }
    & & \\ \\

    \theutterance \stepcounter{utterance}  
    & & & \multicolumn{2}{p{0.3\linewidth}}{
        \cellcolor[rgb]{0.9,0.9,0.9}{
            \makecell[{{p{\linewidth}}}]{
                \texttt{\tiny{[GM$|$GM]}}
                \texttt{Invalid response: RUCH: W, (1, 3)} \\
            }
        }
    }
    & & \\ \\

    \theutterance \stepcounter{utterance}  
    & & \multicolumn{4}{p{0.6\linewidth}}{
        \cellcolor[rgb]{0.9,0.9,0.9}{
            \makecell[{{p{\linewidth}}}]{
                \texttt{\tiny{[P1$\langle$GM]}}
                \texttt{Kara: Musisz rozpocząć grę wysyłając wiadomość do drugiego gracza!} \\
\\ 
\texttt{Upewnij się, że twoja odpowiedź zawiera tylko `POWIEDZ: <WIADOMOŚĆ>` lub `RUCH: <OBIEKT>, (<X>, <Y>)` i nic więcej!} \\
\\ 
\texttt{Łącznie zgromadziliście 2 z 12 kar.} \\
\texttt{Spróbuj ponownie!} \\
            }
        }
    }
    & \\ \\

    \theutterance \stepcounter{utterance}  
    & \multicolumn{4}{p{0.6\linewidth}}{
        \cellcolor[rgb]{0.9,0.9,0.9}{
            \makecell[{{p{\linewidth}}}]{
                \texttt{\tiny{[P1$\rangle$GM]}}
                \texttt{POWIEDZ: Zacznijmy współpracę, aby ustalić wspólną konfigurację obiektów i ukończyć grę przed 20 rundami!} \\
            }
        }
    }
    & & \\ \\

    \theutterance \stepcounter{utterance}  
    & & & \multicolumn{2}{p{0.3\linewidth}}{
        \cellcolor[rgb]{0.9,0.9,0.9}{
            \makecell[{{p{\linewidth}}}]{
                \texttt{\tiny{[GM$|$GM]}}
                \texttt{POWIEDZ: Zacznijmy współpracę, aby ustalić wspólną konfigurację obiektów i ukończyć grę przed 20 rundami!} \\
            }
        }
    }
    & & \\ \\

    \theutterance \stepcounter{utterance}  
    & & \multicolumn{4}{p{0.6\linewidth}}{
        \cellcolor[rgb]{0.9,0.9,0.9}{
            \makecell[{{p{\linewidth}}}]{
                \texttt{\tiny{[GM$\rangle$P2]}}
                \texttt{Jestem mistrzem gry, a Państwo grają w grę kooperacyjną z następującą siatką jako planszą:} \\
\\ 
\texttt{```} \\
\texttt{ 1234567} \\
\texttt{╔═╤══╤══╗} \\
\texttt{║◌│◌◌│◌H║ 1} \\
\texttt{║◌└┬─┴──╢ 2} \\
\texttt{║◌◌│◌◌◌◌║ 3} \\
\texttt{╟─┬┘◌I┌─╢ 4} \\
\texttt{║◌│WT◌│◌║ 5} \\
\texttt{╟─┼─┬─┤◌║ 6} \\
\texttt{║◌│C│◌│◌║ 7} \\
\texttt{╚═╧═╧═╧═╝} \\
\\ 
\texttt{```} \\
\\ 
\texttt{* Górna krawędź wyświetla współrzędne x rosnące w prawo, a prawa krawędź współrzędne y rosnące w dół.} \\
\texttt{* Następujące obiekty są losowo umieszczone na siatce: 'W', 'I', 'T', 'C', 'H'.} \\
\\ 
\\ 
\texttt{Drugi gracz widzi wariant planszy, na której obiekty są umieszczone w różnych losowych miejscach. Nie można zobaczyć planszy drugiego gracza, a on nie może zobaczyć Państwa planszy.} \\
\\ 
\\ 
\texttt{**Cel:**} \\
\\ 
\texttt{Obaj gracze muszą przesuwać obiekty na swoim tle tak, aby identyczne obiekty znalazły się na tych samych współrzędnych. Muszą Państwo komunikować się z drugim graczem, aby uzgodnić wspólną konfigurację docelową.} \\
\\ 
\\ 
\texttt{**Zasady:**} \\
\\ 
\texttt{* W każdej turze można wysłać dokładnie jedno z dwóch poniższych poleceń:} \\
\texttt{ 1. `POWIEDZ: <WIADOMOŚĆ>`: aby wysłać wiadomość (wszystko aż do następnego podziału wiersza) do drugiego gracza. Przekażę ją Państwa partnerowi.} \\
\texttt{2. `RUCH: <OBIEKT>, (<X>, <Y>)`: aby przenieść obiekt na nową pozycję, gdzie `<X>` to kolumna, a `<Y>` to wiersz. Poinformuję Państwa, czy przeniesienie się powiodło, czy nie.} \\
\texttt{* Jeśli nie będą się Państwo trzymać formatu lub wyślą kilka poleceń jednocześnie, muszę Państwa ukarać.} \\
\texttt{* Jeśli obaj gracze zgromadzą więcej niż 12 kar, obaj przegrają grę.} \\
\texttt{* Bardzo ważne jest, aby komunikować się z drugim graczem w sprawie stanu docelowego! *Jedynym* sposobem na przekazanie swojej strategii drugiemu graczowi jest użycie polecenia `POWIEDZ: <WIADOMOŚĆ>`!} \\
\\ 
\\ 
\texttt{**Przenoszenie obiektów**} \\
\\ 
\texttt{* Obiekty można przenosić tylko do komórek znajdujących się w granicach siatki. Komórka docelowa musi być pusta, tzn. musi zawierać tylko symbol '◌'.} \\
\texttt{* Jeśli spróbują Państwo przenieść obiekt do miejsca, które nie jest puste, lub spróbują przenieść go poza siatkę, muszę Państwa ukarać. Otrzymają Państwo kolejną próbę.} \\
\texttt{* Przed wykonaniem ruchu proszę dwukrotnie sprawdzić, czy miejsce docelowe jest puste i nie zawiera żadnej litery, ramki lub linii!} \\
\\ 
\\ 
\texttt{**Koniec gry**} \\
\\ 
\texttt{Jeśli uważają Państwo, że osiągnęli cel wyrównania wszystkich obiektów, mogą Państwo poprosić drugiego gracza o zakończenie gry, wysyłając `POWIEDZ: zakończono?`. Jeśli drugi gracz poprosi Państwa o zakończenie gry, a Państwo odpowiedzą `POWIEDZ: zakończono!`, gra się zakończy.} \\
\\ 
\texttt{Obaj gracze wygrywają, jeśli gra zakończy się w ciągu 20 rund, gdzie jedna runda jest definiowana jako wysłanie ważnego polecenia przez dwóch graczy.} \\
\\ 
\texttt{**Punktacja:**} \\
\\ 
\texttt{Im bliżej siebie znajdują się identyczne obiekty na obu planszach, tym więcej punktów otrzymują Państwo. Kary zmniejszają liczbę punktów. Czy uda się Państwu pobić rekord?} \\
\\ 
\texttt{Drugi gracz rozpoczął grę wysyłając następującą wiadomość:} \\
\\ 
\texttt{"Zacznijmy współpracę, aby ustalić wspólną konfigurację obiektów i ukończyć grę przed 20 rundami!"} \\
\\ 
\texttt{Jakie jest twoje pierwsze polecenie?} \\
            }
        }
    }
    & \\ \\

    \theutterance \stepcounter{utterance}  
    & & & \multicolumn{4}{p{0.6\linewidth}}{
        \cellcolor[rgb]{0.9,0.9,0.9}{
            \makecell[{{p{\linewidth}}}]{
                \texttt{\tiny{[GM$\langle$P2]}}
                \texttt{RUCH: H, (1, 4)} \\
            }
        }
    }
     \\ \\

    \theutterance \stepcounter{utterance}  
    & & & \multicolumn{2}{p{0.3\linewidth}}{
        \cellcolor[rgb]{0.9,0.9,0.9}{
            \makecell[{{p{\linewidth}}}]{
                \texttt{\tiny{[GM$|$GM]}}
                \texttt{RUCH: H, (1, 4)} \\
            }
        }
    }
    & & \\ \\

    \theutterance \stepcounter{utterance}  
    & & & \multicolumn{2}{p{0.3\linewidth}}{
        \cellcolor[rgb]{0.9,0.9,0.9}{
            \makecell[{{p{\linewidth}}}]{
                \texttt{\tiny{[GM$|$GM]}}
                \texttt{Kara: (1,4) nie jest puste, ale zawiera '─'.} \\
\\ 
\texttt{Łącznie zgromadziliście 3 z 12 kar.} \\
\texttt{Spróbuj ponownie!} \\
            }
        }
    }
    & & \\ \\

    \theutterance \stepcounter{utterance}  
    & & \multicolumn{4}{p{0.6\linewidth}}{
        \cellcolor[rgb]{0.9,0.9,0.9}{
            \makecell[{{p{\linewidth}}}]{
                \texttt{\tiny{[GM$\rangle$P2]}}
                \texttt{Kara: (1,4) nie jest puste, ale zawiera '─'.} \\
\\ 
\texttt{Łącznie zgromadziliście 3 z 12 kar.} \\
\texttt{Spróbuj ponownie!} \\
            }
        }
    }
    & \\ \\

    \theutterance \stepcounter{utterance}  
    & & & \multicolumn{4}{p{0.6\linewidth}}{
        \cellcolor[rgb]{0.9,0.9,0.9}{
            \makecell[{{p{\linewidth}}}]{
                \texttt{\tiny{[GM$\langle$P2]}}
                \texttt{RUCH: H, (1, 1)} \\
            }
        }
    }
     \\ \\

    \theutterance \stepcounter{utterance}  
    & & & \multicolumn{2}{p{0.3\linewidth}}{
        \cellcolor[rgb]{0.9,0.9,0.9}{
            \makecell[{{p{\linewidth}}}]{
                \texttt{\tiny{[GM$|$GM]}}
                \texttt{RUCH: H, (1, 1)} \\
            }
        }
    }
    & & \\ \\

    \theutterance \stepcounter{utterance}  
    & & & \multicolumn{2}{p{0.3\linewidth}}{
        \cellcolor[rgb]{0.9,0.9,0.9}{
            \makecell[{{p{\linewidth}}}]{
                \texttt{\tiny{[GM$|$GM]}}
                \texttt{H moved to (1, 1)} \\
            }
        }
    }
    & & \\ \\

    \theutterance \stepcounter{utterance}  
    & & \multicolumn{4}{p{0.6\linewidth}}{
        \cellcolor[rgb]{0.9,0.9,0.9}{
            \makecell[{{p{\linewidth}}}]{
                \texttt{\tiny{[P1$\langle$GM]}}
                \texttt{Twoja wiadomość została przekazana drugiemu graczowi.} \\
\\ 
\\ 
\texttt{Obecnie grasz rundę 1 z maksymalnie 20.} \\
\texttt{Łącznie zgromadziliście 3 z 12 kar.} \\
\texttt{Drugi gracz przesunął obiekt na swojej siatce.} \\
\\ 
\\ 
\texttt{Jakie jest twoje następne polecenie?} \\
            }
        }
    }
    & \\ \\

    \theutterance \stepcounter{utterance}  
    & \multicolumn{4}{p{0.6\linewidth}}{
        \cellcolor[rgb]{0.9,0.9,0.9}{
            \makecell[{{p{\linewidth}}}]{
                \texttt{\tiny{[P1$\rangle$GM]}}
                \texttt{RUCH: W, (1, 3)} \\
            }
        }
    }
    & & \\ \\

    \theutterance \stepcounter{utterance}  
    & & & \multicolumn{2}{p{0.3\linewidth}}{
        \cellcolor[rgb]{0.9,0.9,0.9}{
            \makecell[{{p{\linewidth}}}]{
                \texttt{\tiny{[GM$|$GM]}}
                \texttt{RUCH: W, (1, 3)} \\
            }
        }
    }
    & & \\ \\

    \theutterance \stepcounter{utterance}  
    & & & \multicolumn{2}{p{0.3\linewidth}}{
        \cellcolor[rgb]{0.9,0.9,0.9}{
            \makecell[{{p{\linewidth}}}]{
                \texttt{\tiny{[GM$|$GM]}}
                \texttt{W moved to (1, 3)} \\
            }
        }
    }
    & & \\ \\

    \theutterance \stepcounter{utterance}  
    & & \multicolumn{4}{p{0.6\linewidth}}{
        \cellcolor[rgb]{0.9,0.9,0.9}{
            \makecell[{{p{\linewidth}}}]{
                \texttt{\tiny{[GM$\rangle$P2]}}
                \texttt{Pomyślnie przeniesiono 'H' do (1,1). Twoja zaktualizowana siatka wygląda następująco:} \\
\texttt{```} \\
\texttt{ 1234567} \\
\texttt{╔═╤══╤══╗} \\
\texttt{║H│◌◌│◌◌║ 1} \\
\texttt{║◌└┬─┴──╢ 2} \\
\texttt{║◌◌│◌◌◌◌║ 3} \\
\texttt{╟─┬┘◌I┌─╢ 4} \\
\texttt{║◌│WT◌│◌║ 5} \\
\texttt{╟─┼─┬─┤◌║ 6} \\
\texttt{║◌│C│◌│◌║ 7} \\
\texttt{╚═╧═╧═╧═╝} \\
\texttt{```} \\
\texttt{Obecnie grasz rundę 2 z maksymalnie 20.} \\
\texttt{Łącznie zgromadziliście 3 z 12 kar.} \\
\texttt{Drugi gracz przesunął obiekt na swojej siatce.} \\
\\ 
\\ 
\texttt{Jakie jest twoje następne polecenie?} \\
            }
        }
    }
    & \\ \\

    \theutterance \stepcounter{utterance}  
    & & & \multicolumn{4}{p{0.6\linewidth}}{
        \cellcolor[rgb]{0.9,0.9,0.9}{
            \makecell[{{p{\linewidth}}}]{
                \texttt{\tiny{[GM$\langle$P2]}}
                \texttt{POWIEDZ: Proszę potwierdzić, gdzie znajduje się obiekt 'I' na planszy drugiego gracza, aby móc ustalić wspólną konfigurację docelową.} \\
            }
        }
    }
     \\ \\

    \theutterance \stepcounter{utterance}  
    & & & \multicolumn{2}{p{0.3\linewidth}}{
        \cellcolor[rgb]{0.9,0.9,0.9}{
            \makecell[{{p{\linewidth}}}]{
                \texttt{\tiny{[GM$|$GM]}}
                \texttt{POWIEDZ: Proszę potwierdzić, gdzie znajduje się obiekt 'I' na planszy drugiego gracza, aby móc ustalić wspólną konfigurację docelową.} \\
            }
        }
    }
    & & \\ \\

    \theutterance \stepcounter{utterance}  
    & & \multicolumn{4}{p{0.6\linewidth}}{
        \cellcolor[rgb]{0.9,0.9,0.9}{
            \makecell[{{p{\linewidth}}}]{
                \texttt{\tiny{[P1$\langle$GM]}}
                \texttt{Pomyślnie przeniesiono 'W' do (1,3). Twoja zaktualizowana siatka wygląda następująco:} \\
\texttt{```} \\
\texttt{ 1234567} \\
\texttt{╔═╤══╤══╗} \\
\texttt{║◌│◌◌│◌◌║ 1} \\
\texttt{║◌└┬─┴──╢ 2} \\
\texttt{║W◌│◌◌◌◌║ 3} \\
\texttt{╟─┬┘◌◌┌─╢ 4} \\
\texttt{║I│H◌◌│◌║ 5} \\
\texttt{╟─┼─┬─┤◌║ 6} \\
\texttt{║◌│C│T│◌║ 7} \\
\texttt{╚═╧═╧═╧═╝} \\
\texttt{```} \\
\texttt{Obecnie grasz rundę 2 z maksymalnie 20.} \\
\texttt{Łącznie zgromadziliście 3 z 12 kar.} \\
\texttt{Drugi gracz wysłał tę wiadomość:} \\
\\ 
\\ 
\texttt{"Proszę potwierdzić, gdzie znajduje się obiekt 'I' na planszy drugiego gracza, aby móc ustalić wspólną konfigurację docelową."} \\
\\ 
\\ 
\texttt{Jakie jest twoje następne polecenie?} \\
            }
        }
    }
    & \\ \\

    \theutterance \stepcounter{utterance}  
    & \multicolumn{4}{p{0.6\linewidth}}{
        \cellcolor[rgb]{0.9,0.9,0.9}{
            \makecell[{{p{\linewidth}}}]{
                \texttt{\tiny{[P1$\rangle$GM]}}
                \texttt{POWIEDZ: Obiekt 'I' znajduje się na pozycji 5,2 na mojej planszy.} \\
            }
        }
    }
    & & \\ \\

    \theutterance \stepcounter{utterance}  
    & & & \multicolumn{2}{p{0.3\linewidth}}{
        \cellcolor[rgb]{0.9,0.9,0.9}{
            \makecell[{{p{\linewidth}}}]{
                \texttt{\tiny{[GM$|$GM]}}
                \texttt{POWIEDZ: Obiekt 'I' znajduje się na pozycji 5,2 na mojej planszy.} \\
            }
        }
    }
    & & \\ \\

    \theutterance \stepcounter{utterance}  
    & & \multicolumn{4}{p{0.6\linewidth}}{
        \cellcolor[rgb]{0.9,0.9,0.9}{
            \makecell[{{p{\linewidth}}}]{
                \texttt{\tiny{[GM$\rangle$P2]}}
                \texttt{Twoja wiadomość została przekazana drugiemu graczowi.} \\
\\ 
\\ 
\texttt{Obecnie grasz rundę 3 z maksymalnie 20.} \\
\texttt{Łącznie zgromadziliście 3 z 12 kar.} \\
\texttt{Drugi gracz wysłał tę wiadomość:} \\
\\ 
\\ 
\texttt{"Obiekt 'I' znajduje się na pozycji 5,2 na mojej planszy."} \\
\\ 
\\ 
\texttt{Jakie jest twoje następne polecenie?} \\
            }
        }
    }
    & \\ \\

    \theutterance \stepcounter{utterance}  
    & & & \multicolumn{4}{p{0.6\linewidth}}{
        \cellcolor[rgb]{0.9,0.9,0.9}{
            \makecell[{{p{\linewidth}}}]{
                \texttt{\tiny{[GM$\langle$P2]}}
                \texttt{RUCH: I, (5, 2)} \\
            }
        }
    }
     \\ \\

    \theutterance \stepcounter{utterance}  
    & & & \multicolumn{2}{p{0.3\linewidth}}{
        \cellcolor[rgb]{0.9,0.9,0.9}{
            \makecell[{{p{\linewidth}}}]{
                \texttt{\tiny{[GM$|$GM]}}
                \texttt{RUCH: I, (5, 2)} \\
            }
        }
    }
    & & \\ \\

    \theutterance \stepcounter{utterance}  
    & & & \multicolumn{2}{p{0.3\linewidth}}{
        \cellcolor[rgb]{0.9,0.9,0.9}{
            \makecell[{{p{\linewidth}}}]{
                \texttt{\tiny{[GM$|$GM]}}
                \texttt{Kara: (5,2) nie jest puste, ale zawiera '┴'.} \\
\\ 
\texttt{Łącznie zgromadziliście 4 z 12 kar.} \\
\texttt{Spróbuj ponownie!} \\
            }
        }
    }
    & & \\ \\

    \theutterance \stepcounter{utterance}  
    & & \multicolumn{4}{p{0.6\linewidth}}{
        \cellcolor[rgb]{0.9,0.9,0.9}{
            \makecell[{{p{\linewidth}}}]{
                \texttt{\tiny{[GM$\rangle$P2]}}
                \texttt{Kara: (5,2) nie jest puste, ale zawiera '┴'.} \\
\\ 
\texttt{Łącznie zgromadziliście 4 z 12 kar.} \\
\texttt{Spróbuj ponownie!} \\
            }
        }
    }
    & \\ \\

    \theutterance \stepcounter{utterance}  
    & & & \multicolumn{4}{p{0.6\linewidth}}{
        \cellcolor[rgb]{0.9,0.9,0.9}{
            \makecell[{{p{\linewidth}}}]{
                \texttt{\tiny{[GM$\langle$P2]}}
                \texttt{RUCH: I, (5, 1)} \\
            }
        }
    }
     \\ \\

    \theutterance \stepcounter{utterance}  
    & & & \multicolumn{2}{p{0.3\linewidth}}{
        \cellcolor[rgb]{0.9,0.9,0.9}{
            \makecell[{{p{\linewidth}}}]{
                \texttt{\tiny{[GM$|$GM]}}
                \texttt{RUCH: I, (5, 1)} \\
            }
        }
    }
    & & \\ \\

    \theutterance \stepcounter{utterance}  
    & & & \multicolumn{2}{p{0.3\linewidth}}{
        \cellcolor[rgb]{0.9,0.9,0.9}{
            \makecell[{{p{\linewidth}}}]{
                \texttt{\tiny{[GM$|$GM]}}
                \texttt{Kara: (5,1) nie jest puste, ale zawiera '│'.} \\
\\ 
\texttt{Łącznie zgromadziliście 5 z 12 kar.} \\
\texttt{Spróbuj ponownie!} \\
            }
        }
    }
    & & \\ \\

    \theutterance \stepcounter{utterance}  
    & & \multicolumn{4}{p{0.6\linewidth}}{
        \cellcolor[rgb]{0.9,0.9,0.9}{
            \makecell[{{p{\linewidth}}}]{
                \texttt{\tiny{[GM$\rangle$P2]}}
                \texttt{Kara: (5,1) nie jest puste, ale zawiera '│'.} \\
\\ 
\texttt{Łącznie zgromadziliście 5 z 12 kar.} \\
\texttt{Spróbuj ponownie!} \\
            }
        }
    }
    & \\ \\

    \theutterance \stepcounter{utterance}  
    & & & \multicolumn{4}{p{0.6\linewidth}}{
        \cellcolor[rgb]{0.9,0.9,0.9}{
            \makecell[{{p{\linewidth}}}]{
                \texttt{\tiny{[GM$\langle$P2]}}
                \texttt{RUCH: I, (4, 2)} \\
            }
        }
    }
     \\ \\

    \theutterance \stepcounter{utterance}  
    & & & \multicolumn{2}{p{0.3\linewidth}}{
        \cellcolor[rgb]{0.9,0.9,0.9}{
            \makecell[{{p{\linewidth}}}]{
                \texttt{\tiny{[GM$|$GM]}}
                \texttt{RUCH: I, (4, 2)} \\
            }
        }
    }
    & & \\ \\

    \theutterance \stepcounter{utterance}  
    & & & \multicolumn{2}{p{0.3\linewidth}}{
        \cellcolor[rgb]{0.9,0.9,0.9}{
            \makecell[{{p{\linewidth}}}]{
                \texttt{\tiny{[GM$|$GM]}}
                \texttt{Kara: (4,2) nie jest puste, ale zawiera '─'.} \\
\\ 
\texttt{Łącznie zgromadziliście 6 z 12 kar.} \\
\texttt{Spróbuj ponownie!} \\
            }
        }
    }
    & & \\ \\

    \theutterance \stepcounter{utterance}  
    & & \multicolumn{4}{p{0.6\linewidth}}{
        \cellcolor[rgb]{0.9,0.9,0.9}{
            \makecell[{{p{\linewidth}}}]{
                \texttt{\tiny{[GM$\rangle$P2]}}
                \texttt{Kara: (4,2) nie jest puste, ale zawiera '─'.} \\
\\ 
\texttt{Łącznie zgromadziliście 6 z 12 kar.} \\
\texttt{Spróbuj ponownie!} \\
            }
        }
    }
    & \\ \\

    \theutterance \stepcounter{utterance}  
    & & & \multicolumn{4}{p{0.6\linewidth}}{
        \cellcolor[rgb]{0.9,0.9,0.9}{
            \makecell[{{p{\linewidth}}}]{
                \texttt{\tiny{[GM$\langle$P2]}}
                \texttt{RUCH: I, (5, 3)} \\
            }
        }
    }
     \\ \\

    \theutterance \stepcounter{utterance}  
    & & & \multicolumn{2}{p{0.3\linewidth}}{
        \cellcolor[rgb]{0.9,0.9,0.9}{
            \makecell[{{p{\linewidth}}}]{
                \texttt{\tiny{[GM$|$GM]}}
                \texttt{RUCH: I, (5, 3)} \\
            }
        }
    }
    & & \\ \\

    \theutterance \stepcounter{utterance}  
    & & & \multicolumn{2}{p{0.3\linewidth}}{
        \cellcolor[rgb]{0.9,0.9,0.9}{
            \makecell[{{p{\linewidth}}}]{
                \texttt{\tiny{[GM$|$GM]}}
                \texttt{I moved to (5, 3)} \\
            }
        }
    }
    & & \\ \\

    \theutterance \stepcounter{utterance}  
    & & \multicolumn{4}{p{0.6\linewidth}}{
        \cellcolor[rgb]{0.9,0.9,0.9}{
            \makecell[{{p{\linewidth}}}]{
                \texttt{\tiny{[P1$\langle$GM]}}
                \texttt{Twoja wiadomość została przekazana drugiemu graczowi.} \\
\\ 
\\ 
\texttt{Obecnie grasz rundę 3 z maksymalnie 20.} \\
\texttt{Łącznie zgromadziliście 6 z 12 kar.} \\
\texttt{Drugi gracz przesunął obiekt na swojej siatce.} \\
\\ 
\\ 
\texttt{Jakie jest twoje następne polecenie?} \\
            }
        }
    }
    & \\ \\

    \theutterance \stepcounter{utterance}  
    & \multicolumn{4}{p{0.6\linewidth}}{
        \cellcolor[rgb]{0.9,0.9,0.9}{
            \makecell[{{p{\linewidth}}}]{
                \texttt{\tiny{[P1$\rangle$GM]}}
                \texttt{RUCH: I, (5, 4)} \\
            }
        }
    }
    & & \\ \\

    \theutterance \stepcounter{utterance}  
    & & & \multicolumn{2}{p{0.3\linewidth}}{
        \cellcolor[rgb]{0.9,0.9,0.9}{
            \makecell[{{p{\linewidth}}}]{
                \texttt{\tiny{[GM$|$GM]}}
                \texttt{RUCH: I, (5, 4)} \\
            }
        }
    }
    & & \\ \\

    \theutterance \stepcounter{utterance}  
    & & & \multicolumn{2}{p{0.3\linewidth}}{
        \cellcolor[rgb]{0.9,0.9,0.9}{
            \makecell[{{p{\linewidth}}}]{
                \texttt{\tiny{[GM$|$GM]}}
                \texttt{I moved to (5, 4)} \\
            }
        }
    }
    & & \\ \\

    \theutterance \stepcounter{utterance}  
    & & \multicolumn{4}{p{0.6\linewidth}}{
        \cellcolor[rgb]{0.9,0.9,0.9}{
            \makecell[{{p{\linewidth}}}]{
                \texttt{\tiny{[GM$\rangle$P2]}}
                \texttt{Pomyślnie przeniesiono 'I' do (5,3). Twoja zaktualizowana siatka wygląda następująco:} \\
\texttt{```} \\
\texttt{ 1234567} \\
\texttt{╔═╤══╤══╗} \\
\texttt{║H│◌◌│◌◌║ 1} \\
\texttt{║◌└┬─┴──╢ 2} \\
\texttt{║◌◌│◌I◌◌║ 3} \\
\texttt{╟─┬┘◌◌┌─╢ 4} \\
\texttt{║◌│WT◌│◌║ 5} \\
\texttt{╟─┼─┬─┤◌║ 6} \\
\texttt{║◌│C│◌│◌║ 7} \\
\texttt{╚═╧═╧═╧═╝} \\
\texttt{```} \\
\texttt{Obecnie grasz rundę 4 z maksymalnie 20.} \\
\texttt{Łącznie zgromadziliście 6 z 12 kar.} \\
\texttt{Drugi gracz przesunął obiekt na swojej siatce.} \\
\\ 
\\ 
\texttt{Jakie jest twoje następne polecenie?} \\
            }
        }
    }
    & \\ \\

    \theutterance \stepcounter{utterance}  
    & & & \multicolumn{4}{p{0.6\linewidth}}{
        \cellcolor[rgb]{0.9,0.9,0.9}{
            \makecell[{{p{\linewidth}}}]{
                \texttt{\tiny{[GM$\langle$P2]}}
                \texttt{POWIEDZ: Proszę potwierdzić, gdzie znajduje się obiekt 'W' na planszy drugiego gracza, aby móc ustalić dokładną konfigurację docelową i skończyć grę w 20 rundach.} \\
            }
        }
    }
     \\ \\

    \theutterance \stepcounter{utterance}  
    & & & \multicolumn{2}{p{0.3\linewidth}}{
        \cellcolor[rgb]{0.9,0.9,0.9}{
            \makecell[{{p{\linewidth}}}]{
                \texttt{\tiny{[GM$|$GM]}}
                \texttt{POWIEDZ: Proszę potwierdzić, gdzie znajduje się obiekt 'W' na planszy drugiego gracza, aby móc ustalić dokładną konfigurację docelową i skończyć grę w 20 rundach.} \\
            }
        }
    }
    & & \\ \\

    \theutterance \stepcounter{utterance}  
    & & \multicolumn{4}{p{0.6\linewidth}}{
        \cellcolor[rgb]{0.9,0.9,0.9}{
            \makecell[{{p{\linewidth}}}]{
                \texttt{\tiny{[P1$\langle$GM]}}
                \texttt{Pomyślnie przeniesiono 'I' do (5,4). Twoja zaktualizowana siatka wygląda następująco:} \\
\texttt{```} \\
\texttt{ 1234567} \\
\texttt{╔═╤══╤══╗} \\
\texttt{║◌│◌◌│◌◌║ 1} \\
\texttt{║◌└┬─┴──╢ 2} \\
\texttt{║W◌│◌◌◌◌║ 3} \\
\texttt{╟─┬┘◌I┌─╢ 4} \\
\texttt{║◌│H◌◌│◌║ 5} \\
\texttt{╟─┼─┬─┤◌║ 6} \\
\texttt{║◌│C│T│◌║ 7} \\
\texttt{╚═╧═╧═╧═╝} \\
\texttt{```} \\
\texttt{Obecnie grasz rundę 4 z maksymalnie 20.} \\
\texttt{Łącznie zgromadziliście 6 z 12 kar.} \\
\texttt{Drugi gracz wysłał tę wiadomość:} \\
\\ 
\\ 
\texttt{"Proszę potwierdzić, gdzie znajduje się obiekt 'W' na planszy drugiego gracza, aby móc ustalić dokładną konfigurację docelową i skończyć grę w 20 rundach."} \\
\\ 
\\ 
\texttt{Jakie jest twoje następne polecenie?} \\
            }
        }
    }
    & \\ \\

    \theutterance \stepcounter{utterance}  
    & \multicolumn{4}{p{0.6\linewidth}}{
        \cellcolor[rgb]{0.9,0.9,0.9}{
            \makecell[{{p{\linewidth}}}]{
                \texttt{\tiny{[P1$\rangle$GM]}}
                \texttt{POWIEDZ: Obiekt 'W' znajduje się na pozycji 2,1 na mojej planszy.} \\
            }
        }
    }
    & & \\ \\

    \theutterance \stepcounter{utterance}  
    & & & \multicolumn{2}{p{0.3\linewidth}}{
        \cellcolor[rgb]{0.9,0.9,0.9}{
            \makecell[{{p{\linewidth}}}]{
                \texttt{\tiny{[GM$|$GM]}}
                \texttt{POWIEDZ: Obiekt 'W' znajduje się na pozycji 2,1 na mojej planszy.} \\
            }
        }
    }
    & & \\ \\

    \theutterance \stepcounter{utterance}  
    & & \multicolumn{4}{p{0.6\linewidth}}{
        \cellcolor[rgb]{0.9,0.9,0.9}{
            \makecell[{{p{\linewidth}}}]{
                \texttt{\tiny{[GM$\rangle$P2]}}
                \texttt{Twoja wiadomość została przekazana drugiemu graczowi.} \\
\\ 
\\ 
\texttt{Obecnie grasz rundę 5 z maksymalnie 20.} \\
\texttt{Łącznie zgromadziliście 6 z 12 kar.} \\
\texttt{Drugi gracz wysłał tę wiadomość:} \\
\\ 
\\ 
\texttt{"Obiekt 'W' znajduje się na pozycji 2,1 na mojej planszy."} \\
\\ 
\\ 
\texttt{Jakie jest twoje następne polecenie?} \\
            }
        }
    }
    & \\ \\

    \theutterance \stepcounter{utterance}  
    & & & \multicolumn{4}{p{0.6\linewidth}}{
        \cellcolor[rgb]{0.9,0.9,0.9}{
            \makecell[{{p{\linewidth}}}]{
                \texttt{\tiny{[GM$\langle$P2]}}
                \texttt{RUCH: W, (2, 1)} \\
            }
        }
    }
     \\ \\

    \theutterance \stepcounter{utterance}  
    & & & \multicolumn{2}{p{0.3\linewidth}}{
        \cellcolor[rgb]{0.9,0.9,0.9}{
            \makecell[{{p{\linewidth}}}]{
                \texttt{\tiny{[GM$|$GM]}}
                \texttt{RUCH: W, (2, 1)} \\
            }
        }
    }
    & & \\ \\

    \theutterance \stepcounter{utterance}  
    & & & \multicolumn{2}{p{0.3\linewidth}}{
        \cellcolor[rgb]{0.9,0.9,0.9}{
            \makecell[{{p{\linewidth}}}]{
                \texttt{\tiny{[GM$|$GM]}}
                \texttt{Kara: (2,1) nie jest puste, ale zawiera '│'.} \\
\\ 
\texttt{Łącznie zgromadziliście 7 z 12 kar.} \\
\texttt{Spróbuj ponownie!} \\
            }
        }
    }
    & & \\ \\

    \theutterance \stepcounter{utterance}  
    & & \multicolumn{4}{p{0.6\linewidth}}{
        \cellcolor[rgb]{0.9,0.9,0.9}{
            \makecell[{{p{\linewidth}}}]{
                \texttt{\tiny{[GM$\rangle$P2]}}
                \texttt{Kara: (2,1) nie jest puste, ale zawiera '│'.} \\
\\ 
\texttt{Łącznie zgromadziliście 7 z 12 kar.} \\
\texttt{Spróbuj ponownie!} \\
            }
        }
    }
    & \\ \\

    \theutterance \stepcounter{utterance}  
    & & & \multicolumn{4}{p{0.6\linewidth}}{
        \cellcolor[rgb]{0.9,0.9,0.9}{
            \makecell[{{p{\linewidth}}}]{
                \texttt{\tiny{[GM$\langle$P2]}}
                \texttt{RUCH: W, (2, 5)} \\
            }
        }
    }
     \\ \\

    \theutterance \stepcounter{utterance}  
    & & & \multicolumn{2}{p{0.3\linewidth}}{
        \cellcolor[rgb]{0.9,0.9,0.9}{
            \makecell[{{p{\linewidth}}}]{
                \texttt{\tiny{[GM$|$GM]}}
                \texttt{RUCH: W, (2, 5)} \\
            }
        }
    }
    & & \\ \\

    \theutterance \stepcounter{utterance}  
    & & & \multicolumn{2}{p{0.3\linewidth}}{
        \cellcolor[rgb]{0.9,0.9,0.9}{
            \makecell[{{p{\linewidth}}}]{
                \texttt{\tiny{[GM$|$GM]}}
                \texttt{Kara: (2,5) nie jest puste, ale zawiera '│'.} \\
\\ 
\texttt{Łącznie zgromadziliście 8 z 12 kar.} \\
\texttt{Spróbuj ponownie!} \\
            }
        }
    }
    & & \\ \\

    \theutterance \stepcounter{utterance}  
    & & \multicolumn{4}{p{0.6\linewidth}}{
        \cellcolor[rgb]{0.9,0.9,0.9}{
            \makecell[{{p{\linewidth}}}]{
                \texttt{\tiny{[GM$\rangle$P2]}}
                \texttt{Kara: (2,5) nie jest puste, ale zawiera '│'.} \\
\\ 
\texttt{Łącznie zgromadziliście 8 z 12 kar.} \\
\texttt{Spróbuj ponownie!} \\
            }
        }
    }
    & \\ \\

    \theutterance \stepcounter{utterance}  
    & & & \multicolumn{4}{p{0.6\linewidth}}{
        \cellcolor[rgb]{0.9,0.9,0.9}{
            \makecell[{{p{\linewidth}}}]{
                \texttt{\tiny{[GM$\langle$P2]}}
                \texttt{RUCH: W, (2, 2)} \\
            }
        }
    }
     \\ \\

    \theutterance \stepcounter{utterance}  
    & & & \multicolumn{2}{p{0.3\linewidth}}{
        \cellcolor[rgb]{0.9,0.9,0.9}{
            \makecell[{{p{\linewidth}}}]{
                \texttt{\tiny{[GM$|$GM]}}
                \texttt{RUCH: W, (2, 2)} \\
            }
        }
    }
    & & \\ \\

    \theutterance \stepcounter{utterance}  
    & & & \multicolumn{2}{p{0.3\linewidth}}{
        \cellcolor[rgb]{0.9,0.9,0.9}{
            \makecell[{{p{\linewidth}}}]{
                \texttt{\tiny{[GM$|$GM]}}
                \texttt{Kara: (2,2) nie jest puste, ale zawiera '└'.} \\
\\ 
\texttt{Łącznie zgromadziliście 9 z 12 kar.} \\
\texttt{Spróbuj ponownie!} \\
            }
        }
    }
    & & \\ \\

    \theutterance \stepcounter{utterance}  
    & & \multicolumn{4}{p{0.6\linewidth}}{
        \cellcolor[rgb]{0.9,0.9,0.9}{
            \makecell[{{p{\linewidth}}}]{
                \texttt{\tiny{[GM$\rangle$P2]}}
                \texttt{Kara: (2,2) nie jest puste, ale zawiera '└'.} \\
\\ 
\texttt{Łącznie zgromadziliście 9 z 12 kar.} \\
\texttt{Spróbuj ponownie!} \\
            }
        }
    }
    & \\ \\

    \theutterance \stepcounter{utterance}  
    & & & \multicolumn{4}{p{0.6\linewidth}}{
        \cellcolor[rgb]{0.9,0.9,0.9}{
            \makecell[{{p{\linewidth}}}]{
                \texttt{\tiny{[GM$\langle$P2]}}
                \texttt{RUCH: W, (3, 1)} \\
            }
        }
    }
     \\ \\

    \theutterance \stepcounter{utterance}  
    & & & \multicolumn{2}{p{0.3\linewidth}}{
        \cellcolor[rgb]{0.9,0.9,0.9}{
            \makecell[{{p{\linewidth}}}]{
                \texttt{\tiny{[GM$|$GM]}}
                \texttt{RUCH: W, (3, 1)} \\
            }
        }
    }
    & & \\ \\

    \theutterance \stepcounter{utterance}  
    & & & \multicolumn{2}{p{0.3\linewidth}}{
        \cellcolor[rgb]{0.9,0.9,0.9}{
            \makecell[{{p{\linewidth}}}]{
                \texttt{\tiny{[GM$|$GM]}}
                \texttt{W moved to (3, 1)} \\
            }
        }
    }
    & & \\ \\

    \theutterance \stepcounter{utterance}  
    & & \multicolumn{4}{p{0.6\linewidth}}{
        \cellcolor[rgb]{0.9,0.9,0.9}{
            \makecell[{{p{\linewidth}}}]{
                \texttt{\tiny{[P1$\langle$GM]}}
                \texttt{Twoja wiadomość została przekazana drugiemu graczowi.} \\
\\ 
\\ 
\texttt{Obecnie grasz rundę 5 z maksymalnie 20.} \\
\texttt{Łącznie zgromadziliście 9 z 12 kar.} \\
\texttt{Drugi gracz przesunął obiekt na swojej siatce.} \\
\\ 
\\ 
\texttt{Jakie jest twoje następne polecenie?} \\
            }
        }
    }
    & \\ \\

    \theutterance \stepcounter{utterance}  
    & \multicolumn{4}{p{0.6\linewidth}}{
        \cellcolor[rgb]{0.9,0.9,0.9}{
            \makecell[{{p{\linewidth}}}]{
                \texttt{\tiny{[P1$\rangle$GM]}}
                \texttt{RUCH: W, (2, 4)} \\
            }
        }
    }
    & & \\ \\

    \theutterance \stepcounter{utterance}  
    & & & \multicolumn{2}{p{0.3\linewidth}}{
        \cellcolor[rgb]{0.9,0.9,0.9}{
            \makecell[{{p{\linewidth}}}]{
                \texttt{\tiny{[GM$|$GM]}}
                \texttt{RUCH: W, (2, 4)} \\
            }
        }
    }
    & & \\ \\

    \theutterance \stepcounter{utterance}  
    & & & \multicolumn{2}{p{0.3\linewidth}}{
        \cellcolor[rgb]{0.9,0.9,0.9}{
            \makecell[{{p{\linewidth}}}]{
                \texttt{\tiny{[GM$|$GM]}}
                \texttt{Kara: (2,4) nie jest puste, ale zawiera '┬'.} \\
\\ 
\texttt{Łącznie zgromadziliście 10 z 12 kar.} \\
\texttt{Spróbuj ponownie!} \\
            }
        }
    }
    & & \\ \\

    \theutterance \stepcounter{utterance}  
    & & \multicolumn{4}{p{0.6\linewidth}}{
        \cellcolor[rgb]{0.9,0.9,0.9}{
            \makecell[{{p{\linewidth}}}]{
                \texttt{\tiny{[P1$\langle$GM]}}
                \texttt{Kara: (2,4) nie jest puste, ale zawiera '┬'.} \\
\\ 
\texttt{Łącznie zgromadziliście 10 z 12 kar.} \\
\texttt{Spróbuj ponownie!} \\
            }
        }
    }
    & \\ \\

    \theutterance \stepcounter{utterance}  
    & \multicolumn{4}{p{0.6\linewidth}}{
        \cellcolor[rgb]{0.9,0.9,0.9}{
            \makecell[{{p{\linewidth}}}]{
                \texttt{\tiny{[P1$\rangle$GM]}}
                \texttt{RUCH: W, (2, 2)} \\
            }
        }
    }
    & & \\ \\

    \theutterance \stepcounter{utterance}  
    & & & \multicolumn{2}{p{0.3\linewidth}}{
        \cellcolor[rgb]{0.9,0.9,0.9}{
            \makecell[{{p{\linewidth}}}]{
                \texttt{\tiny{[GM$|$GM]}}
                \texttt{RUCH: W, (2, 2)} \\
            }
        }
    }
    & & \\ \\

    \theutterance \stepcounter{utterance}  
    & & & \multicolumn{2}{p{0.3\linewidth}}{
        \cellcolor[rgb]{0.9,0.9,0.9}{
            \makecell[{{p{\linewidth}}}]{
                \texttt{\tiny{[GM$|$GM]}}
                \texttt{Kara: (2,2) nie jest puste, ale zawiera '└'.} \\
\\ 
\texttt{Łącznie zgromadziliście 11 z 12 kar.} \\
\texttt{Spróbuj ponownie!} \\
            }
        }
    }
    & & \\ \\

    \theutterance \stepcounter{utterance}  
    & & \multicolumn{4}{p{0.6\linewidth}}{
        \cellcolor[rgb]{0.9,0.9,0.9}{
            \makecell[{{p{\linewidth}}}]{
                \texttt{\tiny{[P1$\langle$GM]}}
                \texttt{Kara: (2,2) nie jest puste, ale zawiera '└'.} \\
\\ 
\texttt{Łącznie zgromadziliście 11 z 12 kar.} \\
\texttt{Spróbuj ponownie!} \\
            }
        }
    }
    & \\ \\

    \theutterance \stepcounter{utterance}  
    & \multicolumn{4}{p{0.6\linewidth}}{
        \cellcolor[rgb]{0.9,0.9,0.9}{
            \makecell[{{p{\linewidth}}}]{
                \texttt{\tiny{[P1$\rangle$GM]}}
                \texttt{RUCH: W, (0, 3)} \\
            }
        }
    }
    & & \\ \\

    \theutterance \stepcounter{utterance}  
    & & & \multicolumn{2}{p{0.3\linewidth}}{
        \cellcolor[rgb]{0.9,0.9,0.9}{
            \makecell[{{p{\linewidth}}}]{
                \texttt{\tiny{[GM$|$GM]}}
                \texttt{RUCH: W, (0, 3)} \\
            }
        }
    }
    & & \\ \\

    \theutterance \stepcounter{utterance}  
    & & & \multicolumn{2}{p{0.3\linewidth}}{
        \cellcolor[rgb]{0.9,0.9,0.9}{
            \makecell[{{p{\linewidth}}}]{
                \texttt{\tiny{[GM$|$GM]}}
                \texttt{Kara: (0,3) nie jest puste, ale zawiera '║'.} \\
\\ 
\texttt{Łącznie zgromadziliście 12 z 12 kar.} \\
\texttt{Spróbuj ponownie!} \\
            }
        }
    }
    & & \\ \\

    \theutterance \stepcounter{utterance}  
    & & \multicolumn{4}{p{0.6\linewidth}}{
        \cellcolor[rgb]{0.9,0.9,0.9}{
            \makecell[{{p{\linewidth}}}]{
                \texttt{\tiny{[P1$\langle$GM]}}
                \texttt{Kara: (0,3) nie jest puste, ale zawiera '║'.} \\
\\ 
\texttt{Łącznie zgromadziliście 12 z 12 kar.} \\
\texttt{Spróbuj ponownie!} \\
            }
        }
    }
    & \\ \\

    \theutterance \stepcounter{utterance}  
    & \multicolumn{4}{p{0.6\linewidth}}{
        \cellcolor[rgb]{0.9,0.9,0.9}{
            \makecell[{{p{\linewidth}}}]{
                \texttt{\tiny{[P1$\rangle$GM]}}
                \texttt{Nie mogę kontynuować tej sesji. Proszę zacząć od nowa, podając aktualną planszę lub polecenie zgodne z zasadami.} \\
            }
        }
    }
    & & \\ \\

    \theutterance \stepcounter{utterance}  
    & & & \multicolumn{2}{p{0.3\linewidth}}{
        \cellcolor[rgb]{0.9,0.9,0.9}{
            \makecell[{{p{\linewidth}}}]{
                \texttt{\tiny{[GM$|$GM]}}
                \texttt{Nie mogę kontynuować tej sesji. Proszę zacząć od nowa, podając aktualną planszę lub polecenie zgodne z zasadami.} \\
            }
        }
    }
    & & \\ \\

    \theutterance \stepcounter{utterance}  
    & & & \multicolumn{2}{p{0.3\linewidth}}{
        \cellcolor[rgb]{0.9,0.9,0.9}{
            \makecell[{{p{\linewidth}}}]{
                \texttt{\tiny{[GM$|$GM]}}
                \texttt{Invalid response format} \\
            }
        }
    }
    & & \\ \\

    \theutterance \stepcounter{utterance}  
    & & & \multicolumn{2}{p{0.3\linewidth}}{
        \cellcolor[rgb]{0.9,0.9,0.9}{
            \makecell[{{p{\linewidth}}}]{
                \texttt{\tiny{[GM$|$GM]}}
                \texttt{Maximum number of penalties exceeded} \\
            }
        }
    }
    & & \\ \\

    \theutterance \stepcounter{utterance}  
    & & & \multicolumn{2}{p{0.3\linewidth}}{
        \cellcolor[rgb]{0.9,0.9,0.9}{
            \makecell[{{p{\linewidth}}}]{
                \texttt{\tiny{[GM$|$GM]}}
                \texttt{Initial states:} \\
\texttt{```} \\
\texttt{Player 1:} \\
\texttt{ 1234567} \\
\texttt{╔═╤══╤══╗} \\
\texttt{║◌│◌◌│◌◌║ 1} \\
\texttt{║W└┬─┴──╢ 2} \\
\texttt{║◌◌│◌◌◌◌║ 3} \\
\texttt{╟─┬┘◌◌┌─╢ 4} \\
\texttt{║I│H◌◌│◌║ 5} \\
\texttt{╟─┼─┬─┤◌║ 6} \\
\texttt{║◌│C│T│◌║ 7} \\
\texttt{╚═╧═╧═╧═╝} \\
\\ 
\\ 
\texttt{Player 2:} \\
\texttt{ 1234567} \\
\texttt{╔═╤══╤══╗} \\
\texttt{║◌│◌◌│◌H║ 1} \\
\texttt{║◌└┬─┴──╢ 2} \\
\texttt{║◌◌│◌◌◌◌║ 3} \\
\texttt{╟─┬┘◌I┌─╢ 4} \\
\texttt{║◌│WT◌│◌║ 5} \\
\texttt{╟─┼─┬─┤◌║ 6} \\
\texttt{║◌│C│◌│◌║ 7} \\
\texttt{╚═╧═╧═╧═╝} \\
\\ 
\texttt{```} \\
            }
        }
    }
    & & \\ \\

    \theutterance \stepcounter{utterance}  
    & & & \multicolumn{2}{p{0.3\linewidth}}{
        \cellcolor[rgb]{0.9,0.9,0.9}{
            \makecell[{{p{\linewidth}}}]{
                \texttt{\tiny{[GM$|$GM]}}
                \texttt{End states:} \\
\texttt{```} \\
\texttt{Player 1:} \\
\texttt{ 1234567} \\
\texttt{╔═╤══╤══╗} \\
\texttt{║◌│◌◌│◌◌║ 1} \\
\texttt{║◌└┬─┴──╢ 2} \\
\texttt{║W◌│◌◌◌◌║ 3} \\
\texttt{╟─┬┘◌I┌─╢ 4} \\
\texttt{║◌│H◌◌│◌║ 5} \\
\texttt{╟─┼─┬─┤◌║ 6} \\
\texttt{║◌│C│T│◌║ 7} \\
\texttt{╚═╧═╧═╧═╝} \\
\\ 
\\ 
\texttt{Player 2:} \\
\texttt{ 1234567} \\
\texttt{╔═╤══╤══╗} \\
\texttt{║H│W◌│◌◌║ 1} \\
\texttt{║◌└┬─┴──╢ 2} \\
\texttt{║◌◌│◌I◌◌║ 3} \\
\texttt{╟─┬┘◌◌┌─╢ 4} \\
\texttt{║◌│◌T◌│◌║ 5} \\
\texttt{╟─┼─┬─┤◌║ 6} \\
\texttt{║◌│C│◌│◌║ 7} \\
\texttt{╚═╧═╧═╧═╝} \\
\\ 
\texttt{```} \\
            }
        }
    }
    & & \\ \\

    \theutterance \stepcounter{utterance}  
    & & & \multicolumn{2}{p{0.3\linewidth}}{
        \cellcolor[rgb]{0.9,0.9,0.9}{
            \makecell[{{p{\linewidth}}}]{
                \texttt{\tiny{[GM$|$GM]}}
                \texttt{* success: False} \\
\texttt{* lose: False} \\
\texttt{* aborted: True} \\
\texttt{{-}{-}{-}{-}{-}{-}{-}} \\
\texttt{* Shifts: 3.00} \\
\texttt{* Max Shifts: 8.00} \\
\texttt{* Min Shifts: 4.00} \\
\texttt{* End Distance Sum: 10.54} \\
\texttt{* Init Distance Sum: 15.62} \\
\texttt{* Expected Distance Sum: 20.95} \\
\texttt{* Penalties: 13.00} \\
\texttt{* Max Penalties: 12.00} \\
\texttt{* Rounds: 5.00} \\
\texttt{* Max Rounds: 20.00} \\
\texttt{* Object Count: 5.00} \\
            }
        }
    }
    & & \\ \\

\end{supertabular}
}

\end{document}
